\section{USAJMO 2024}
        \begin{enumerate}
        	\item (\textit{USAJMO 2024}) Let $ABCD$ be a cyclic quadrilateral with $AB = 7$ and $CD = 8$. Points $P$ and $Q$ are selected on segment $AB$ such that $AP = BQ = 3$. Points $R$ and $S$ are selected on segment $CD$ such that $CR = DS = 2$. Prove that $PQRS$ is a cyclic quadrilateral.


 

	\item (\textit{USAJMO 2024}) Let $m$ and $n$ be positive integers. Let $S$ be the set of integer points $(x,y)$ with $1\leq x\leq2m$ and $1\leq y\leq2n$. A configuration of $mn$ rectangles is called \textit{happy} if each point in $S$ is a vertex of exactly one rectangle, and all rectangles have sides parallel to the coordinate axes. Prove that the number of \textit{happy} configurations is odd.


 

	\item (\textit{USAJMO 2024}) Let $a(n)$ be the sequence defined by $a(1)=2$ and $a(n+1)=(a(n))^{n+1}-1$ for each integer $n\geq1$. Suppose that $p>2$ is prime and $k$ is a positive integer. Prove that some term of the sequence $a(n)$ is divisible by $p^k$.


 

	\item (\textit{USAJMO 2024}) Let $n \ge 3$ be an integer. Rowan and Colin play a game on an $n \times n$ grid of squares, where each square is colored either red or blue. Rowan is allowed to permute the rows of the grid, and Colin is allowed to permute the columns of the grid. A grid coloring is $orderly$ if:


 
                \begin{itemize}
                    \item no matter how Rowan permutes the rows of the coloring, Colin can then permute the columns to restore the original grid coloring; and
                \end{itemize}


                
                \begin{itemize}
                    \item no matter how Colin permutes the columns of the coloring, Rowan can then permute the rows to restore the original grid coloring;
                \end{itemize}


                In terms of $n$, how many orderly colorings are there?


 

	\item (\textit{USAJMO 2024}) Find all functions $f:\mathbb{R}\rightarrow\mathbb{R}$ that satisfy

 \[f(x^2-y)+2yf(x)=f(f(x))+f(y)\]
for all $x,y\in\mathbb{R}$.


 

	\item (\textit{USAJMO 2024}) Point $D$ is selected inside acute triangle $ABC$ so that $\angle DAC=\angle ACB$ and $\angle BDC=90^\circ+\angle BAC$. Point $E$ is chosen on ray $BD$ so that $AE=EC$. Let $M$ be the midpoint of $BC$. Show that line $AB$ is tangent to the circumcircle of triangle $BEM$.


 

\end{enumerate}